% =============================================================
% CAPÍTULO 7 — Conclusiones y trabajo futuro
% Versión: matrícula de honor
% =============================================================

\section{Conclusiones y trabajo futuro}
\label{sec:conclusiones}

Este capítulo recapitula las contribuciones del trabajo, discute con
honestidad sus limitaciones y propone líneas concretas de investigación
futura. Frente a la práctica habitual de cerrar memorias con una lista de
hallazgos, optamos por una estructura que distingue tres niveles: lo que el
trabajo describe sobre la organización interna de BERT, lo que el
trabajo diseña como contribución metodológica, y lo que el trabajo
abre como agenda de investigación.


% -------------------------------------------------------------
\subsection{Contribuciones principales}
\label{sec:contribuciones-finales}
% -------------------------------------------------------------

Este trabajo ha producido un mapa funcional de cómo BERT-base
\emph{fine-tuneado} organiza internamente la información emocional sobre
GoEmotions (23~emociones), utilizando la compresión SVD selectiva como
instrumento experimental y triangulando los hallazgos con cinco técnicas
de interpretabilidad mecánica. Las contribuciones se organizan en tres
bloques: la cartografía funcional propiamente dicha, el uso de la
compresión como diagnóstico estructural y la traducción del mapa en una
estrategia de compresión informada.

La Tabla~\ref{tab:objetivos-cumplimiento} resume el cumplimiento de los
12~sub-objetivos planteados en el Anexo de gestión (\S\ref{sec:gep-alcance}),
indicando para cada uno la sección de la memoria donde se desarrolla el
resultado y la evidencia cuantitativa más relevante. De los 11~objetivos
del núcleo, los 11~se han alcanzado; el O12 (recuperación por
\emph{fine-tuning}, originalmente opcional) se promovió a tarea principal
tras observar el efecto regularizador documentado en
\S\ref{sec:finetuning-recovery}.

\begin{table}[ht]
\centering
\caption{Cumplimiento de los sub-objetivos planteados en el plan inicial.
Sección indica dónde se desarrolla el resultado; evidencia, la métrica
clave que documenta el cumplimiento.}
\label{tab:objetivos-cumplimiento}
\footnotesize
\begin{tabular}{@{}p{0.45cm} p{4.6cm} p{2.9cm} p{4.6cm}@{}}
\toprule
\textbf{Obj.} & \textbf{Descripción} & \textbf{Sección} & \textbf{Evidencia} \\
\midrule
O1  & Baseline competitivo (BERT 23 emociones) & \S\ref{sec:met-modelo}     & F1 macro 0{,}577 sobre test \\
O2  & Módulo SVD reutilizable                 & \S\ref{sec:met-svd}        & \texttt{src/compression/svd.py} \\
O3  & Trade-off compresión--rendimiento       & \S\ref{sec:compresion-uniforme} & Cliff entre $r{=}384$ y $r{=}256$ \\
O4  & Visualización degradación               & \S\ref{sec:probing}        & Fig.~\ref{fig:cap5-galaxia-evolution} \\
O5  & Sensibilidad por componente             & \S\ref{sec:sensibilidad-componente} & Asimetría 14$\times$ a $r{=}128$ \\
O6  & Sensibilidad por profundidad            & \S\ref{sec:sensibilidad-componente} & Late catastrófica a $r{=}128$ \\
O7  & Compresión adaptativa por energía       & \S\ref{sec:compresion-adaptativa} & Adaptive 99\%: 96{,}7\% retención \\
O8  & Localizar cristalización emocional      & \S\ref{sec:probing}        & 276 probes; 9/8/6 emociones por banda \\
O9  & Evidencia funcional de localización     & \S\ref{sec:logit-lens}, \S\ref{sec:activation-patching} & FFN L11 sola = 100\% restauración \\
O10 & Mapear cabezas y neuronas               & \S\ref{sec:head-ablation}, \S\ref{sec:neuron-analysis} & 144 cabezas + 36\,864 neuronas \\
O11 & Compresión informada                    & \S\ref{sec:greedy-algorithm} & 8/9 puntos Pareto-óptimos greedy \\
O12 & Recuperación por \emph{fine-tuning}      & \S\ref{sec:finetuning-recovery} & F1 macro $\geq$ baseline con 86{,}4\% params (preliminar) \\
\bottomrule
\end{tabular}
\end{table}

\subsubsection{Cartografía funcional del modelo}

La aplicación coordinada de cinco técnicas de interpretabilidad (probing
lineal, \emph{logit lens}, \emph{activation patching}, ablación de cabezas y análisis de
selectividad neuronal) ha producido un mapa funcional coherente de cómo
BERT codifica y procesa la información emocional. Las observaciones se
articulan en cuatro hallazgos:

\begin{itemize}
  \item Arquitectura funcional de tres niveles. Las capas tempranas
        (Emb--L2) extraen señal léxica bruta y capturan aproximadamente el
        61\% de la separabilidad lineal final del modelo
        ($F1_{\text{probe}}^{L0} = 0{,}349$ frente al máximo
        $F1_{\text{probe}}^{L11} = 0{,}569$, \S\ref{sec:probing}). Las
        capas intermedias (L3--L7) realizan cómputo de transición en una
        geometría desalineada con la base del clasificador, lo que se
        manifiesta en el valle del \emph{logit lens} (\S\ref{sec:logit-lens}). Las
        capas profundas (L8--L11) ejecutan la cristalización: alinean la
        representación con la cabeza de clasificación entrenada y suprimen
        activamente las opciones incorrectas. Este último paso constituye el
        cuello de botella geométrico del modelo y explica su criticidad
        bajo compresión.

  \item Cristalización emocional progresiva determinada por la
        semántica, no por la frecuencia. Las 23 emociones se distribuyen en
        tres bandas de cristalización con 9, 8 y 6 emociones respectivamente.
        La frecuencia en el dataset no predice la profundidad de
        cristalización (Spearman $\rho = -0{,}26$, $p = 0{,}22$,
        \S\ref{sec:probing}), mientras que la profundidad sí predice
        fuertemente el rendimiento alcanzable
        ($\rho = -0{,}67$, $p = 0{,}0004$). Emociones con señales léxicas
        claras (\emph{gratitude}, \emph{love}, \emph{amusement}) cristalizan
        en L0; emociones que requieren desambiguación pragmática
        (\emph{annoyance}, \emph{disappointment}) cristalizan en L8--L9. Esta
        observación tiene implicaciones que van más allá de este trabajo:
        sugiere que añadir datos a las emociones contextualmente complejas no
        adelanta su cristalización.

  \item La FFN de la capa~11 como locus computacional crítico.
        Restaurar exclusivamente este componente (aproximadamente un tercio
        de los parámetros de una capa) en un modelo previamente colapsado
        por compresión recupera el 100\% del rendimiento baseline
        (\S\ref{sec:activation-patching}). Este hallazgo desafía
        productivamente la narrativa centrada en la atención que ha
        prevalecido desde \emph{Attention Is All You Need}
        \cite{vaswani2017attention}. La interpretación mecanicista, articulada
        a través del \emph{logit lens}, es que la FFN de las capas tardías realiza
        la \emph{rotación geométrica} que lleva la representación
        contextualizada a la base sobre la que opera la cabeza de
        clasificación.

  \item Gradiente monotónico de criticidad atencional. La proporción
        de cabezas críticas crece de forma estricta con la profundidad: 27\%
        en capas 0--4, 53\% en capas 5--7, 77\% en capas 8--11. La capa 11
        constituye un caso extremo: ninguna de sus 12 cabezas es prescindible
        (\S\ref{sec:head-ablation}). Identificamos 21~cabezas con
        \emph{efecto interferente} (su ablación mejora el rendimiento), que
        constituyen candidatas directas a poda sin re-entrenamiento.
\end{itemize}


\subsubsection{Compresión como instrumento de diagnóstico}

El análisis sistemático de compresión SVD, más allá de su finalidad aplicada,
ha funcionado como un instrumento de diagnóstico que revela la organización
interna del modelo de un modo complementario a las técnicas de
interpretabilidad mecánica:

\begin{itemize}
  \item Asimetría radical entre componentes. La compresión a
        rango~128 produce retenciones de F1 que difieren en un factor 14$\times$
        entre el extremo más robusto (Query: 99{,}4\%) y el más frágil (FFN
        Intermediate: 6{,}9\%). En el régimen severo de rango~64 la asimetría
        se amplía hasta el factor 72$\times$ normalizado por parámetros
        eliminados: Query mantiene una retención del 98{,}6\%, mientras que
        FFN Intermediate colapsa al 0\% (\S\ref{sec:sensibilidad-componente}). Esta asimetría es consistente
        con la estructura espectral de los pesos: las matrices Q y K presentan
        decaimiento espectral pronunciado, mientras que las FFN exhiben un
        espectro plano.

  \item Transición de fase bajo compresión uniforme. La compresión
        uniforme presenta un \emph{cliff} abrupto entre rango~384 (43{,}5\%
        de retención) y rango~256 (4{,}3\%). A rangos inferiores el modelo
        entra en colapso clínico (F1 = 0). La existencia de esta transición
        de fase, junto con la asimetría entre componentes, demuestra que la
        compresión uniforme es estructuralmente incapaz de explorar la
        región interesante del espacio compresión--rendimiento.

  \item Predictores de la vulnerabilidad emocional. Tres
        propiedades del modelo predicen significativamente la caída de F1
        bajo compresión: la profundidad de cristalización por probing
        ($\rho \approx 0{,}48$), el aislamiento representacional
        ($\rho = 0{,}45$) y, como predictor más fuerte, la norma del
        vector de selectividad neuronal ($\rho = 0{,}64$, $p = 0{,}001$,
        \S\ref{sec:neuron-analysis}). Las emociones cuyo modelo escribe con
        mayor intensidad neuronal son las más vulnerables a perturbaciones
        generales de los pesos, aunque la SVD no las dañe selectivamente.

  \item Disociación geometría SVD--función. La fracción de señal
        emocional que cae en el subespacio eliminado por SVD es
        prácticamente idéntica para las 23 emociones (rango 0{,}53--0{,}56),
        y su correlación con la caída de F1 es nula
        ($\rho = -0{,}004$, $p = 0{,}99$). La SVD opera en un espacio
        matemático que no ``ve'' la función computacional del modelo: elimina
        señal uniformemente. Este resultado, aparentemente negativo, es
        teóricamente importante: la geometría que la norma de Frobenius
        optimiza es ortogonal a la geometría de la especialización
        neuronal.
\end{itemize}


\subsubsection{Compresión informada por interpretabilidad}

La conexión entre los dos ejes del trabajo se materializa en una estrategia
de compresión data-driven que constituye, en términos prácticos, la
contribución más directamente aplicable de la memoria:

\begin{itemize}
  \item Dominancia de la frontera de Pareto. Un algoritmo greedy
        alimentado con datos empíricos de sensibilidad por componente y
        profundidad produce 8 de los 9 puntos Pareto-óptimos sobre 21
        configuraciones evaluadas (6 uniformes, 4 adaptativas, 3 heurísticas
        informadas y 8 greedy). Al ratio de parámetros del 80\%, la
        estrategia greedy retiene el 87\% del F1 macro, frente al 43\% de
        la compresión uniforme y el 25\% de la adaptativa energética
        (\S\ref{sec:greedy-algorithm}).

  \item Validación cruzada por convergencia algorítmica. La
        secuencia de decisiones del algoritmo greedy, que opera
        exclusivamente con datos numéricos de sensibilidad, reproduce
        independientemente los hallazgos cualitativos de interpretabilidad:
        comprime primero Q y K (casi gratuitas), después FFN Output en
        capas tempranas, y nunca toca la FFN Intermediate de las capas
        tardías (\S\ref{sec:greedy-algorithm}). Esta convergencia, entre dos
        líneas metodológicamente independientes, refuerza la validez de
        ambas: si interpretabilidad y compresión llegan a las mismas
        conclusiones por caminos disjuntos, la probabilidad de que ambas
        estén sistemáticamente sesgadas en la misma dirección disminuye.

  \item Resultado negativo informativo: la heurística no basta.
        Las tres estrategias informadas basadas en reglas cualitativas
        (\textit{informed-light}, \textit{informed-moderate},
        \textit{informed-aggressive}) convergen exactamente sobre la
        frontera de las baselines ciegas: \texttt{informed\_aggressive}
        reproduce \texttt{uniform\_r256}, \texttt{informed\_moderate}
        reproduce \texttt{uniform\_r512}, e \texttt{informed\_light}
        requiere un 28{,}5\% más parámetros que el original
        (\S\ref{sec:informed-heuristic}). Ninguna domina la frontera de
        Pareto. Este resultado delimita con precisión el valor aplicado
        de la interpretabilidad mecánica: el conocimiento cualitativo
        identifica qué investigar, pero solo la medición empírica de
        sensibilidad permite determinar cuánto comprimir. La diferencia
        entre saber que la FFN es más importante que la atención y saber
        que Q a rango~256 retiene 100{,}4\% del F1 mientras que FFN
        Intermediate a rango~64 colapsa al 0\% es la diferencia entre
        intuición y diseño.

  \item Recuperación por fine-tuning y observación de regularización
        implícita. Un modelo comprimido al 86{,}4\% de los parámetros, tras
        tres épocas de \emph{fine-tuning}, alcanza un F1 macro superior al baseline
        sin comprimir. Las ganancias se concentran sistemáticamente en las
        emociones infrarrepresentadas, con \emph{embarrassment} pasando de
        F1 = 0{,}267 a F1 = 0{,}509 ($+90{,}6\%$ relativo,
        \S\ref{sec:finetuning-recovery}). Este patrón es consistente con la
        hipótesis de que la SVD actúa como regularización estructural,
        al eliminar direcciones de baja energía que pueden codificar
        sobreajuste a patrones espurios. Su verificación
        causal requeriría experimentos de control que se proponen en
        \S\ref{sec:trabajo-futuro}.
\end{itemize}


% -------------------------------------------------------------
\subsection{Limitaciones}
\label{sec:limitaciones}
% -------------------------------------------------------------

Los resultados de este trabajo deben interpretarse dentro de un alcance
experimental delimitado, cuyas restricciones se enumeran a continuación con
honestidad metodológica.

\subsubsection{Generalización: un modelo, una tarea}

Todos los experimentos se realizan sobre un único modelo (BERT-base, 110\,M
parámetros) y una única tarea (clasificación emocional multi-etiqueta sobre
GoEmotions). No es posible afirmar que los patrones observados (la
jerarquía de sensibilidad Q/K $\ll$ V $\ll$ FFN, la concentración funcional
en capas tardías, la correlación cristalización-vulnerabilidad, la
dominancia de la FFN sobre la atención en \emph{activation patching}) se
mantengan en otros modelos o en otras tareas. La arquitectura encoder-only de
BERT, con su cabeza de clasificación sobre el token \texttt{[CLS]}, podría
favorecer la concentración de información en capas finales de un modo que no
se reproduzca en arquitecturas decoder-only o encoder-decoder. La extensión
sistemática a BERT-large, RoBERTa, GPT-2 y LLaMA constituye la línea de
trabajo futuro de mayor prioridad (\S\ref{sec:trabajo-futuro}).

\subsubsection{Sesgos del dataset}

GoEmotions presenta desequilibrios notables entre clases. Cinco emociones
fueron excluidas del análisis (neutral, \emph{grief},
\emph{nervousness}, \emph{pride}, \emph{relief}) por dos razones distintas
que conviene no confundir: neutral se excluyó por no constituir una
emoción propiamente dicha y por dominar el dataset (más de 14\,000 ejemplos),
distorsionando las métricas macro-promediadas; las cuatro restantes se
excluyeron por producir F1 = 0 en el baseline tras \emph{fine-tuning}, debido a su
bajísima frecuencia (menos de 100 ejemplos en entrenamiento). Las 23
emociones restantes varían en frecuencia desde 303 ejemplos
(\emph{embarrassment}) hasta 4\,130 (\emph{admiration}), un factor 13$\times$.
Aunque la frecuencia no predice la profundidad de cristalización, sí podría
afectar la robustez estadística de las métricas de especialización neuronal y
la sensibilidad emocional a la compresión, especialmente en emociones con
menos de 500 ejemplos.

\subsubsection{Compresión post-hoc, no compresión durante entrenamiento}

La compresión se ha limitado a SVD aplicada después del \emph{fine-tuning}, sin
re-entrenamiento durante la compresión y sin combinación con otras técnicas.
Otras aproximaciones (poda estructurada, cuantización, destilación de
conocimiento o entrenamiento consciente de la compresión) podrían interactuar de forma diferente con
la estructura funcional del modelo y posiblemente alcanzar mejores
trade-offs. En particular, el pruning estructurado de cabezas de atención
podría beneficiarse directamente de la taxonomía presentada en
\S\ref{sec:head-ablation}, donde las 38 cabezas dispensable y las 21
con efecto interferente son candidatas naturales a eliminación sin
re-entrenamiento.

\subsubsection{Limitaciones de las técnicas de interpretabilidad}

\paragraph{Probes lineales.} Los \textit{probing classifiers} utilizados son
regresiones logísticas con regularización por defecto ($C = 1{,}0$),
deliberadamente simples siguiendo la metodología de Hewitt y
Liang~\cite{hewitt2019designing}. Esta elección tiene ventajas (mide la
geometría de la representación, no la capacidad del clasificador) pero
también limitaciones: información emocional codificada de forma no lineal en
capas tempranas no sería detectada. Es posible que las capas Emb--L2
contengan más información emocional de la que el probing lineal documenta.

\paragraph{\emph{Logit lens} y calibración.} La aplicación del clasificador final a
representaciones de capas intermedias produce predicciones cuyas magnitudes
absolutas no son directamente comparables al baseline, ya que la cabeza de
clasificación se entrenó exclusivamente sobre representaciones de L11 tras
LayerNorm. Reportamos los resultados con esta cautela: las cifras absolutas
funcionan como diagnóstico de tendencias, no como medida calibrada. La
forma cualitativa del patrón en U es robusta y se corrobora con probing y
patching, pero las magnitudes específicas (gold sigmoid de 0{,}02 en L7,
0{,}44 en L11) deben interpretarse en términos relativos. Una versión más
rigurosa con \emph{tuned lens} \cite{belrose2023eliciting} se discute en
\S\ref{sec:trabajo-futuro}.

\paragraph{\emph{Activation patching} y efecto cadena.} El \emph{activation patching}
opera restaurando pesos originales en un modelo previamente comprimido, lo
que introduce un acoplamiento que conviene reconocer: la contribución medida
de una capa depende del estado de las capas no restauradas. Esto explica por
qué las capas tempranas muestran recuperación nula a pesar de ser
informativas según probing: su señal limpia se destruye al atravesar 10
capas comprimidas posteriores. La corrupción mediante SVD truncada a
rango~64 tampoco es neutral en sentido estadístico (a diferencia del
\emph{causal tracing} clásico de Meng et al.\ \cite{meng2022locating},
que usa ruido gaussiano), sino una distorsión estructurada. Las
conclusiones de patching tienen por tanto naturaleza funcional (qué
componentes son suficientes para reactivar el modelo desde el colapso)
más que estrictamente causal en sentido de
Pearl. La distinción es importante para futuras réplicas.

\subsubsection{Potencia estadística}

Varias correlaciones centrales del trabajo (cristalización vs.\
vulnerabilidad, $\rho \approx 0{,}48$; norma de selectividad vs.\ caída de
F1, $\rho = 0{,}64$) se calculan sobre $n = 23$ puntos (uno por emoción).
Aunque las correlaciones más relevantes son altamente significativas
($p \leq 0{,}001$), la potencia estadística es limitada y no permite
detectar efectos moderados con confianza.

Cuantitativamente, mediante la transformación $z$ de Fisher con
$n=23$, $\alpha=0{,}05$ y potencia $1-\beta = 0{,}80$, la mínima
correlación de Spearman detectable es
$|\rho_{\min}| \approx 0{,}556$. Las correlaciones fuertes
documentadas en la memoria (cristalización vs.\ F1 máximo:
$\rho = -0{,}673$; norma de selectividad neuronal vs.\ caída de F1:
$\rho = +0{,}639$) están \emph{justo dentro} del rango detectable, lo
que es consistente con sus $p$-valores por debajo de $0{,}001$. En
contraste, la correlación frecuencia--cristalización ($\rho = -0{,}265$)
no es detectable a nuestro $n$: para detectarla con $\alpha=0{,}05$ y
potencia $0{,}80$ se necesitarían aproximadamente $n \geq 85$ emociones.
La conclusión metodológica es directa: con $n=23$, los efectos
\emph{grandes} ($|\rho| > 0{,}55$) son robustamente detectables, los
efectos \emph{moderados} ($0{,}3 \leq |\rho| \leq 0{,}55$) requieren
muestras mayores, y los efectos \emph{pequeños} ($|\rho| < 0{,}3$) son
inobservables. Los intervalos de confianza bootstrap al 95\% reportados
en \S\ref{sec:probing} reflejan esta limitación con honestidad: la
amplitud de los intervalos para correlaciones moderadas
(p.~ej.\ frecuencia--cristalización $[-0{,}65, +0{,}19]$) cruza cero,
señalando ausencia de evidencia, no evidencia de ausencia.

\subsubsection{Recuperación por \emph{fine-tuning}: ausencia de grupo de control}

La afirmación de que la compresión actúa como regularización implícita se
basa en la observación de que un modelo comprimido + \emph{fine-tuned} supera al
baseline sin comprimir. Sin embargo, no se ha entrenado un grupo de control
consistente en el baseline original con tres épocas adicionales de
\emph{fine-tuning}, lo que impide aislar el efecto de la compresión del efecto del
re-entrenamiento adicional. Es posible, aunque improbable dada la
saturación de F1 macro observada en epochs 3-4, que parte de la mejora
provenga simplemente de más entrenamiento. Este experimento de control se
propone en \S\ref{sec:trabajo-futuro} como verificación causal de la
hipótesis de regularización.


% -------------------------------------------------------------
\subsection{Trabajo futuro}
\label{sec:trabajo-futuro}
% -------------------------------------------------------------

Las limitaciones identificadas y los hallazgos inesperados del trabajo
sugieren seis líneas concretas de investigación futura, ordenadas por
prioridad metodológica.

\subsubsection{Generalización a otros modelos y tareas}

La pregunta más inmediata es si la jerarquía de sensibilidad por componente
y la correlación cristalización--vulnerabilidad se reproducen en modelos de
mayor escala (BERT-large, RoBERTa-large) y en arquitecturas decoder-only
(GPT-2, LLaMA). Un estudio comparativo permitiría determinar qué hallazgos
son propiedades universales de la arquitectura Transformer y cuáles son
específicos de BERT-base o de la tarea emocional. De particular interés sería
verificar si la dominancia funcional de la FFN sobre la atención
documentada en \S\ref{sec:activation-patching} se mantiene en modelos con más
capas y dimensiones ocultas mayores. La hipótesis de Geva et
al.~\cite{geva2021transformer} sobre las FFN como memorias clave-valor
predice que sí, pero requiere verificación empírica directa con la
metodología de este trabajo.

\subsubsection{Verificación causal del efecto regularizador}

La hipótesis de que la compresión SVD funciona como regularización
implícita merece verificación experimental rigurosa. El diseño propuesto
incluiría tres condiciones de comparación: (i) baseline sin compresión + 3
epochs adicionales de \emph{fine-tuning}, (ii) baseline con compresión greedy +
3 epochs de \emph{fine-tuning} (configuración actual), y (iii) baseline + 3 epochs
con dropout/weight decay incrementado (regularización explícita). Si la
condición (ii) supera a (i) y (iii), la hipótesis de regularización
estructural por SVD obtiene apoyo causal. Si (i) iguala a (ii), el efecto
observado se explica por re-entrenamiento adicional. Este experimento es
factible con los recursos disponibles y constituye la verificación más
inmediata pendiente.

\subsubsection{Compresión a granularidad de cabeza individual}

La estrategia greedy actual opera a granularidad de componente y banda de
profundidad. Sin embargo, los resultados de ablación muestran que dentro de
una misma capa coexisten cabezas críticas y dispensables
(\S\ref{sec:head-ablation}). Una extensión natural sería un algoritmo de
compresión que asigne rangos SVD individuales a cada cabeza de atención en
función de su importancia medida, lo que podría mejorar la frontera de
Pareto especialmente en las capas intermedias (L5--L7), donde la
composición de cabezas críticas y dispensables es más heterogénea. La
combinación con la categoría \emph{Dispensable} de la taxonomía
(38 cabezas) eliminada directamente sin SVD constituiría un híbrido
poda-SVD que podría alcanzar ratios de compresión sustancialmente mayores.

\subsubsection{\emph{Tuned lens} para corrección del sesgo de calibración}

La aplicación del \emph{logit lens} que documentamos sufre de un sesgo conocido:
las representaciones de capas intermedias no están calibradas con la cabeza
de clasificación entrenada exclusivamente sobre L11. Belrose et
al.~\cite{belrose2023eliciting} proponen \textit{tuned lens}: entrenar una
transformación lineal $T_l: \mathbb{R}^d \to \mathbb{R}^d$ por capa que
minimice la divergencia KL entre la predicción del clasificador aplicado a
$T_l(h_l)$ y la predicción real del modelo. Aplicar \emph{tuned lens} a este caso
permitiría calibrar el patrón en U y verificar si el valle observado
(L4--L8) es robusto a la corrección de calibración o si se atenúa, lo que
afinaría la interpretación geométrica de las capas medias.

\subsubsection{Combinación de técnicas de compresión}

Los resultados sugieren una estrategia de compresión compuesta que combine
SVD con poda estructurada y cuantización. La taxonomía de cabezas del
\S\ref{sec:head-ablation} identifica 38 cabezas dispensable y 21 con
efecto interferente, candidatas a eliminación directa sin re-entrenamiento.
Un pipeline en cascada \textit{poda-de-cabezas-interferentes
$\rightarrow$ poda-de-cabezas-dispensable $\rightarrow$ SVD greedy
$\rightarrow$ cuantización post-SVD $\rightarrow$ \emph{fine-tuning}} permitiría
componer las reducciones de los cuatro métodos. La pregunta científica
abierta es si las reducciones se acumulan multiplicativamente (cada técnica
explora dimensiones de redundancia ortogonales) o sub-multiplicativamente
(las técnicas compiten por la misma redundancia).

\subsubsection{Interpretabilidad dinámica durante el entrenamiento}

Todos los análisis de este trabajo operan sobre un modelo ya entrenado. Una
extensión ambiciosa consistiría en monitorizar la cristalización y la
especialización neuronal durante el proceso de \emph{fine-tuning}, no solo al
final. Esto permitiría observar cómo emergen los patrones documentados
(cuándo cristaliza cada emoción, cuándo se establece la dominancia de la
FFN, cuándo aparecen las cabezas interferentes) y eventualmente informar
decisiones de \textit{early stopping} o adaptación dinámica del learning
rate por capa. La integración con frameworks de interpretabilidad como
TransformerLens permitiría implementar este análisis sin coste computacional
prohibitivo.

\subsubsection{Validación en entorno de producción}

Los tiempos de inferencia y el consumo de memoria de los modelos
comprimidos no se han medido en este trabajo, que se ha centrado en la
retención de F1 macro. Una evaluación de latencia y throughput en hardware
real (CPU, GPU, dispositivos edge) complementaría el análisis y determinaría
si las ventajas teóricas de la compresión informada se traducen en
beneficios prácticos de despliegue. Esta extensión es relativamente directa y
debería preceder a cualquier propuesta de adopción industrial de la
metodología.


% -------------------------------------------------------------
\subsection{Predicciones falsables del framework}
\label{sec:predicciones-falsables}
% -------------------------------------------------------------

Un trabajo empírico trasciende el caso estudiado cuando genera
predicciones contrastables sobre escenarios fuera del alcance
experimental original. Las cinco predicciones que siguen se derivan
del framework descrito en esta memoria; cualquier réplica futura sobre
otros modelos o tareas puede confirmarlas o refutarlas, delimitando
empíricamente qué hallazgos son específicos de BERT-base sobre
GoEmotions y cuáles son propiedades estructurales más generales.

La primera predicción concierne a la \emph{asimetría espectral entre
Q/K y FFN en BERT-large}. Si esta asimetría
(\S\ref{sec:analisis-espectral}) es una propiedad arquitectónica del
Transformer encoder y no un artefacto de la pequeña escala de
BERT-base, BERT-large debería exhibir un cociente similar entre el
rango efectivo de Q/K y FFN. Cuantitativamente, esperamos
$k_{95}(\text{Q})/k_{95}(\text{FFN}) \in [0{,}55, 0{,}75]$
(BERT-base presenta $0{,}64$). Un cociente cercano a $1$ en BERT-large
indicaría que la asimetría es específica del régimen de baja capacidad
y obligaría a revisar la generalización.

La segunda predicción distingue \emph{encoder-only de decoder-only}.
El hallazgo de \S\ref{sec:activation-patching} (restaurar L11 sola
recupera el 100\%) descansa sobre la geometría de clasificación a
través del token \texttt{[CLS]}. En modelos decoder-only sin cabeza de
clasificación dedicada (GPT-2, LLaMA), el patching debería distribuirse
más uniformemente entre las capas tardías; replicar la misma
concentración extrema en una sola capa obligaría a revisar la
explicación geométrica que articula el logit lens.

La tercera concierne a la \emph{independencia de la cristalización
respecto a la arquitectura}. La distribución 9/8/6 emociones por banda
de cristalización (\S\ref{sec:probing}) y la correlación
cristalización--F1 máximo ($\rho = -0{,}67$) deberían replicarse
cualitativamente en RoBERTa-base sobre GoEmotions, dada la similitud
arquitectónica. Si el orden de cristalización de emociones
\emph{individuales} cambia marcadamente entre modelos, la profundidad
de procesamiento dependería del corpus de pre-entrenamiento
(Wikipedia + BookCorpus para BERT vs.\ corpus extendido para RoBERTa)
y no únicamente de la complejidad lingüística de la emoción.

La cuarta predice que el \emph{efecto regularizador de la SVD se
reproduce en otras tareas con clases desequilibradas}. La observación
de \S\ref{sec:finetuning-recovery} (compresión + \emph{fine-tuning}
beneficia desproporcionadamente a las emociones infrarrepresentadas)
debería aparecer en NER con etiquetas raras o en clasificación
multi-etiqueta sobre datasets desbalanceados como Reuters-21578 o
AAPD. La ausencia del efecto refutaría la hipótesis de que la
regularización implícita procede de la eliminación de direcciones de
baja energía, o exigiría condiciones adicionales no identificadas en
este trabajo.

La quinta es la más útil prácticamente: \emph{la heurística
cualitativa converge sobre la frontera ciega también en otros
dominios}. El resultado negativo de \S\ref{sec:informed-heuristic}
(las tres heurísticas informadas igualan a las baselines uniformes)
predice un fenómeno general: cualquier traducción cualitativa de
hallazgos interpretativos a reglas de compresión convergerá sobre lo
que la compresión ciega alcanza, salvo que se incorporen datos
cuantitativos de sensibilidad. Convertir esta observación en una guía
para futuros trabajos de ``compresión informada por interpretabilidad''
ahorraría meses de trabajo en proyectos similares.

Las cinco predicciones son ortogonales: cada una falla bajo condiciones
distintas. La confirmación conjunta supondría un test estricto del
framework; la refutación individual delimitaría con precisión qué parte
del análisis es generalizable y cuál pertenece al caso particular.


% -------------------------------------------------------------
\subsection{Reflexión final}
\label{sec:reflexion-final}
% -------------------------------------------------------------

Este trabajo partió de una pregunta aparentemente técnica (cuánto se puede
comprimir un modelo de lenguaje sin degradar su rendimiento) y derivó en
una investigación sobre la estructura funcional interna del modelo. El
camino recorrido sugiere una observación metodológica que excede el caso
estudiado: la compresión y la interpretabilidad no son disciplinas
disjuntas, sino dos perspectivas sobre el mismo objeto.

La compresión es interpretabilidad implícita: cuando comprimimos un
componente y observamos qué se rompe, estamos midiendo su contribución
funcional con la herramienta más directa posible: su ausencia. La
interpretabilidad es compresión racional: cuando entendemos qué hace cada
componente, sabemos cuáles podemos eliminar sin perder función. El
resultado negativo de la heurística (\S\ref{sec:informed-heuristic}) y el
éxito del algoritmo greedy (\S\ref{sec:greedy-algorithm}) precisan esta
afirmación: la conexión entre ambas disciplinas no es cualitativa sino
cuantitativa. Los hallazgos de probing, patching y ablación describen la
arquitectura funcional, pero solo los datos numéricos de sensibilidad
permiten traducir esa descripción en una estrategia de compresión efectiva.

El hallazgo central de la memoria, en esta lectura, no es ninguno de los
resultados individuales (la asimetría 72$\times$ entre componentes, la FFN
de la capa 11 como cuello de botella, la dominancia de la frontera de
Pareto, el efecto regularizador del \emph{fine-tuning}), sino su convergencia.
Cinco técnicas de interpretabilidad, dos familias de compresión y un
algoritmo de optimización ciego apuntan al mismo eje funcional: las capas
8--11, y especialmente sus FFN, son donde reside la capacidad emocional del
modelo. Que cinco metodologías independientes lleguen a la misma conclusión
hace la conclusión más robusta que cualquiera de ellas individualmente.

Una segunda observación, más especulativa, concierne al carácter de los
modelos de lenguaje pre-entrenados como objetos de estudio. BERT no fue
diseñado para clasificar emociones; emergió de la competencia general
adquirida durante el pre-entrenamiento masked-language-modeling sobre Wikipedia
y BookCorpus. El \emph{fine-tuning} sobre GoEmotions no creó la arquitectura
funcional documentada aquí: la organizó. Las capas tardías como cuello de
botella geométrico, la dominancia de la FFN, la cristalización progresiva,
el patrón en U del \emph{logit lens}: ninguno de estos fenómenos fue explícitamente
optimizado durante el entrenamiento. Emergieron. Esto sugiere que las
estructuras funcionales de los Transformers fine-tuneados son, en sentido
fuerte, \emph{descubrimientos empíricos}, no propiedades de diseño. La
interpretabilidad mecánica, vista así, es la disciplina que documenta lo
que la optimización por gradiente decidió, no lo que sus diseñadores
prescribieron.

Si este trabajo cumple su función, debería contribuir modestamente a esa
disciplina, ofreciendo un mapa anatómico parcial pero coherente de cómo un
Transformer organiza el procesamiento emocional, y demostrando que ese
mapa puede informar la compresión más allá de lo que las heurísticas
permiten. Que el resultado se obtenga sobre un modelo y una tarea concretos
limita su alcance; que la metodología sea replicable amplifica su utilidad.
Las extensiones a otros modelos y tareas, propuestas en
\S\ref{sec:trabajo-futuro}, son la prueba pendiente.